% Source: https://artofproblemsolving.com/community/c3h2635211_easy_geometry_for_you_to_try
% Stable identity: aops-topic-2635211-q01
Since $\angle ABC=90^\circ$, the hypotenuse $AC$ is a diameter of the circumcircle, so the circumcentre $O$ is the midpoint of $\overline{AC}$. By construction $O$ is also the midpoint of $\overline{BB'}$.

\begin{asy}
usepackage("amsmath");
size(4.5cm);
defaultpen(fontsize(10pt));
pair B=(0,0),A=(0,3),C=(4,0),O=(2,1.5),Bp=(4,3);
draw(circle(O,2.5),gray);
draw(A--B--C--cycle,linewidth(0.8));draw(A--Bp--C,linewidth(0.8));
draw(B--Bp,dashed);draw(A--C,dashed);
draw((0,0.3)--(0.3,0.3)--(0.3,0));
dot(O);label("$O$",O,N);
label("$A$",A,NW);label("$B$",B,SW);label("$C$",C,SE);label("$B'$",Bp,NE);
\end{asy}

So the diagonals $AC$ and $BB'$ of quadrilateral $ABCB'$ bisect each other, and $ABCB'$ is a parallelogram. A parallelogram with a right angle is a rectangle (consecutive angles of a parallelogram are supplementary, and opposite angles are equal, so all four angles are $90^\circ$). Since $\angle ABC=90^\circ$, $ABCB'$ is a rectangle. (Also $BB'=2R=AC$: the diagonals are equal, as they must be in a rectangle.) $\blacksquare$
